\clearpage\fancyhead[L]{\small INSTRUMENTE.}\hypertarget{native-370-lower}{}\label{native-370-lower}
\begin{center}\textbf{INSTRUMENTE.}\par\vspace{4pt}\hrule\vspace{4pt}\fontsize{10}{12}\selectfont Göttingische gelehrte Anzeigen. Stück 75. Seite 745\;..\;752. 1813. Mai 10.\par\vspace{4pt}\hrule\vspace{4pt}\end{center}

\par\indent In der Sitzung der königl. Societät der Wissenschaften am 3. April, in welcher der Professor \textsc{Eichhorn} die Vorlesung hielt, ertheilte der Professor \textsc{Gauss} eine Nachricht von zwei neuen astronomischen Instrumenten, welche die hiesige Sternwarte aus der Werkstatt des berühmten Künstlers Hrn. Salinenrath \textsc{Reichenbach} in München seit einigen Monaten besitzt, und legte zugleich die Erstlinge der damit angestellten Beobachtungen vor.

\par\indent Das \emph{erstere} ist ein Repetitionskreis mit zwei Fernröhren, von 12 Zoll Durchmesser. Man bewundert an diesem Kunstwerke eben so sehr die Feinheit und Genauigkeit der Theilung, die fast unglaubliche Empfindlichkeit der Libellen, die Vollkommenheit der Fernröhre, als die Accuratesse und Schönheit der Ausarbeitung aller einzelnen Theile des Instruments, die Leichtigkeit aller der vielfältigen Bewegungen und die mancherlei sinnreich angeordneten Vorrichtungen, wodurch der geniale Künstler für die Bequemlichkeit des Beobachters gesorgt hat.

\par\indent Der Hauptkreis ist auf eingelegtem Silber unmittelbar von 5 zu 5 Minuten getheilt, und jeder der vier Nonien theilt einen solchen Theil wiederum in 75 Theile, also von 4 zu 4 Secunden. Diese Nonien, unter rechten Winkeln von einander abstehend, befinden sich, in derselben Ebene mit dem Gradbogen, auf einem vollständigen zweiten Kreise, der sich innerhalb des erstern so vollkommen concentrisch bewegt, dass er, obgleich mit blossen Augen gar kein Zwischenraum wahrgenommen wird, \SourceLeaf{371}doch jenen nirgends berührt, und daher die Bewegung mit grösster Leichtigkeit von statten
geht. Auch durch die absichtlich etwas schief gestellten Microscope bemerkt man diesen
Zwischenraum nicht, sondern die Striche der Nonien scheinen unmittelbar an die Striche des
Gradbogens zu stossen. Dadurch wird die Genauigkeit des Ablesens sehr befördert, so wie
durch eine zweckmässige Beleuchtung, und durch ungemeine Zartheit und Gleichheit der
Theilstriche, und man kann füglich 2, ja allenfalls einzelne Secunden, schätzen. Die vier
Nonien weichen nie mehr als ein paar Secunden von einander ab, wodurch sich sowohl die
Abwesenheit aller merklichen Excentricität, als die unübertreffliche Genauigkeit der
Eintheilung selbst beweist.

\begingroup\rightskip=0pt plus 1em \par\indent Die drei zum Kreise gehörigen Libellen sind von ausserordentlicher Empfindlichkeit.
Das Haupt-Niveau am hintern Fernrohr gibt auf eine Secunde Neigung einen Ausschlag von mehr
als Einer Pariser Linie. Diese so äusserst geringe und doch gleichförmige Krümmung im
Innern der Glasröhre, welche einen Halbmesser von mehr als 1400 Fuss voraussetzt, konnte nur
durch eine sehr künstliche und delicate Bearbeitung erhalten werden, und die Glasröhren sind
daher, um diese Krümmung nicht wieder zu verlieren, an ihren beiden Enden nicht
zugeschmolzen, sondern mit genau eingeschliffenen Glasstöpseln, worüber noch eine
Blasenhaut gezogen ist, auf das vollkommenste verschlossen.\par\endgroup

\par\indent Die beiden Fernröhre sind, obgleich sie nur eine Länge von 16 Zoll und eine
Oeffnung von 15 Linien haben, doch von ganz ausgezeichneter Güte, aufs vollkommenste
achromatisch, und vertragen bei ihrer ungemeinen Präcision starke Vergrösserungen. Das
Pointiren auf den feinen, im Brennpunkt eingezogenen, Spinnefäden geschieht daher mit
grosser Schärfe. Das vordere Fernrohr hat ein prismatisches Ocular, so dass beim
Höhenmessen das Auge immer durch horizontale Strahlen sieht, und daher hohe Sterne sich
eben so bequem, wie niedrige, beobachten lassen. Der Verlust an Licht ist dabei fast ganz
unmerklich.

\par\indent Alle die Schrauben, welche zur feinern Stellung dienen, sind mit grösster
Accuratesse gearbeitet. Auf die leiseste Berührung sprechen sie sogleich gehörig an, und
durch besondere Bremsschrauben ist auch auf die Zukunft allem todten Gange vorgebeugt.
Ihre Feinheit ist so gross, dass über 100 Gänge auf die Länge eines Zolls gehen.

\par\indent Wir übergehen, da hier nicht der Ort zu einer vollständigen Beschreibung des
Instruments ist, mehrere eben so sinnreiche als zweckmässige Vorrichtungen, welche den
\textsc{Reichenbach}schen Kreisen eigenthümlich sind, z.\,B. die Mittel, die Ebene des
Kreises auf das genaueste vertical zu stellen, die Gegenstände mit Leichtigkeit
aufzufinden und in das Gesichtsfeld zu bringen und dergl.

\par\indent Mehrere Umstände verzögerten den Anfang der astronomischen Beobachtungen mit
diesem Kreise bis zur Mitte des März. Von den Beobachtungen, welche seitdem Prof.
\textsc{Gauss}, mit Beihülfe des Hrn. Prof. \textsc{Harding}, der die Einstellung des
Niveaus besorgte, gemacht hat, theilen wir hier als Probe die Resultate für die Polhöhe
mit, die sich aus den beobachteten untern Culminationen des Polarsterns ergeben haben,
und deren schöne Uebereinstimmung am besten für die Vortrefflichkeit des Instruments
zeugt.

\Needspace{13\baselineskip}\begin{center}Polhöhe der Göttinger Sternwarte aus Beobachtungen des Polarsterns, in der untern Columination\end{center}
\begin{center}\fontsize{10}{12}\selectfont\setlength{\tabcolsep}{7pt}\renewcommand{\arraystretch}{1.15}\begin{tabular}{ll|c|r}
& 1813 & \shortstack{Anzahl der\\Beobachtungen} &\\\hline
März & 20 & 10 & $51^\circ\;31'\;54''86$\\
& 22 & 18 & $55.73$\\
& 26 & 18 & $57.41$\\
& 31 & 18 & $56.25$\\
April & 3 & 18 & $56.55$\\
& 7 & 32 & $55.20$\\
& 8 & 22 & $57.33$\\
\end{tabular}\end{center}

\par\indent\SourceLeaf{372}Das Mittel aus allen 136 Beobachtungen ist $51^\circ\,31'\,56''20$, wovon das äusserste Resultat nur $1''34$ abweicht. Hievon geht noch ab $0''16$ Reduction auf den Mittelpunkt der Sternwarte, deren Polhöhe folglich $51^\circ\,31'\,56''04$ wird, nur $2''$ grösser, als sie \textsc{Tob. Meyer} bestimmt hat. Dies Resultat ist noch abhängig von der Declination des Polarsterns, welche nach Hrn. von \textsc{Zach}'s Bestimmung zum Grunde gelegt ist, und wird daher vielleicht, wenn der Polarstern erst in der obern Culmination beobachtet werden kann, noch eine, aber gewiss sehr kleine, Modification erleiden. Hätte man für die Declination das Mittel aus von \textsc{Zach}'s, \textsc{Oriani}'s, \textsc{Bouvard}'s und \textsc{Pond}'s Bestimmungen (deren Extreme nur $\frac{3}{4}$ Secunden von einander abweichen) zum Grunde gelegt, oder \textsc{Pond}'s Bestimmung allein, welche mit diesem Mittel genau übereinstimmt, so wäre die Polhöhe noch um $0''44$ kleiner ausgefallen.

\par\indent Das \emph{zweite} Instrument, in seiner Art ein eben so bewundernswürdiges Meisterwerk, ist ein Repetitionstheodolith von 8 Zoll Durchmesser. Der Horizontalkreis ist unmittelbar von 10 zu 10 Minuten getheilt; jeder der vier Nonien gibt $10''$; kleinere Theile lassen sich noch schätzen. Auch hier sind die Nonien auf einer vollständigen Kreisscheibe innerhalb des getheilten Kreises, und mit diesem in derselben Ebene. Der innere oder Noniuskreis trägt zwei Stützen, auf welchen das Hauptfernrohr von 12 Zoll Länge und 13 Linien Oeffnung, gerade wie ein Passageninstrument, an einer horizontalen, in zwei vollkommen gleiche cylindrische stählerne Zapfen auslaufenden, Axe aufgehängt ist, und dessen Gesichtslinie auch, gerade wie die eines Mittagsfernrohres, durch Umlegen auf das genaueste auf diese Axe senkrecht gebracht werden kann. Um die Ebene des Kreises horizontal, und die erwähnte Axe ihr parallel zu stellen, dient ein an der Axe anzuhängendes Niveau, bei welchem 1 Secunde einen Ausschlag von einer halben Linie gibt. Das Fernrohr kann, ehe es den Horizontalkreis berührt, bis zu $40^\circ$ über und unter den Horizont geneigt werden, und diese Neigung wird an einem besondern Verticalkreise von 5 Zoll Durchmesser gemessen, dessen Nonius einzelne Minuten gibt, und an welchem halbe oder Drittelminuten sich noch schätzen lassen. Das zweite, untere Fernrohr ist dem obern ganz gleich, dient aber nur als Versicherungsfernrohr für den unbeweglichen Stand des Instruments. Es erregt Erstaunen, wie genau sich mit einem so kleinen Instrumente Winkel messen. lassen. Einzelne Messungen geben die Winkel allemal bis auf wenige Secunden genau, und durch Repetition kann man sich der wahren Werthe der Winkel bis auf $1''$, höchstens $2''$, versichern. Vorzüglich wichtig für den astronomischen Gebrauch ist die Leichtigkeit und Genauigkeit, womit man, vermittelst Beobachtung der Sonne oder der Fixsterne (wovon man die von der ersten Grösse mit Leichtigkeit am hellen Tage sieht), Azimuthe irdischer Gegenstände bestimmen kann, so wie diese, wenn sie für einen gewissen Standplatz einmal scharf ausgemittelt sind, wiederum zu bequemen und von der Refraction unabhängigen Zeitbestimmungen dienen können: ein Vortheil, der besonders in den Wintermonaten auf einer Sternwarte, welche kein Passageninstrument besitzt, sehr hoch anzuschlagen ist.

\par\indent Besonders erfreulich ist, dass die schönen Fernröhre an diesen Instrumenten, welche an Vollkommenheit den besten englischen von gleichen Dimensionen nichts nachgeben, sondern sie eher noch übertreffen, ganz deutschen Ursprunges sind: das Flintglas zu den Objectiven wird in Benedictbeuern verfertigt. --- [Dieser Anzeige hat \textsc{Gauss} ein Verzeichniss der Preise der in dem optischen Institute von \textsc{Reichenbach}, \textsc{Frauenhofer} und von \textsc{Utzschneider} verfertigten Sehwerkzeuge beigefügt, das hier in der Ausgab seiner Werke nicht mit aufgenommen ist.]
